\chapter{Input}

\section{Three questions you can ask about a key}

This is short but it matters enormously, because picking the wrong one of the
three is the most common beginner bug in any engine.

\begin{center}
\begin{tabular}{@{}lll@{}}
\toprule
\textbf{Function} & \textbf{True when} & \textbf{Use it for} \\
\midrule
\lstinline|IsKeyDown(KEY_D)| & every frame it is held & walking, aiming, holding \\
\lstinline|IsKeyPressed(KEY_D)| & the single frame it goes down & jump, shoot, menus \\
\lstinline|IsKeyReleased(KEY_D)| & the single frame it comes up & charge-and-release \\
\bottomrule
\end{tabular}
\end{center}

Hold a key for one second at 60 FPS: \lstinline|IsKeyDown| is true sixty times,
\lstinline|IsKeyPressed| is true once.

\gotcha{Use \lstinline|IsKeyDown| for a menu and your cursor jumps sixty entries
in one second. Use \lstinline|IsKeyPressed| for walking and your character
teleports one step per tap and refuses to walk. When input ``feels broken'',
check which of the two you used before anything else.}

Key names are the obvious ones: \lstinline|KEY_A| to \lstinline|KEY_Z|,
\lstinline|KEY_ZERO| to \lstinline|KEY_NINE|, \lstinline|KEY_SPACE|,
\lstinline|KEY_ENTER|, \lstinline|KEY_LEFT_SHIFT|, \lstinline|KEY_LEFT_CONTROL|,
\lstinline|KEY_UP|, \lstinline|KEY_LEFT|, \lstinline|KEY_TAB|,
\lstinline|KEY_ESCAPE|. They are all listed in \file{raylib.h} under
\lstinline|KeyboardKey|.

\gotcha{\lstinline|KEY_ESCAPE| closes the window by default --- that is what
\lstinline|WindowShouldClose()| watches. If you want Escape to open a pause menu
instead, turn the default off with \lstinline|SetExitKey(KEY_NULL)| and handle it
yourself.}

\section{The mouse}

\begin{lstlisting}
Vector2 mouse = GetMousePosition();      // in SCREEN pixels
float   wheel = GetMouseWheelMove();     // -1, 0 or +1 sized steps
Vector2 delta = GetMouseDelta();         // how far it moved THIS frame

if (IsMouseButtonPressed(MOUSE_BUTTON_LEFT))  { /* click */ }
if (IsMouseButtonDown(MOUSE_BUTTON_RIGHT))    { /* held */ }
\end{lstlisting}

The same pressed/down/released trio applies to buttons.

\lstinline|GetMouseDelta()| is the one that makes first person cameras possible.
It reports movement rather than position, which is exactly what you want when
the cursor is locked in place and can no longer move:

\begin{lstlisting}
DisableCursor();     // hide the cursor and lock it to the window centre
EnableCursor();      // give it back to the user
\end{lstlisting}

With the cursor disabled, \lstinline|GetMousePosition()| stops being useful
(the pointer is pinned) but \lstinline|GetMouseDelta()| keeps reporting how much
the user moved the physical mouse. Chapter 12 turns that into a camera.

\gotcha{Always call \lstinline|EnableCursor()| before \lstinline|CloseWindow()|,
or a crash can leave the user's pointer captured and invisible. Every lesson in
the 3D course does this.}

\section{Text input}

For a name entry field you do not want key codes, you want characters, with the
keyboard layout and the shift key already applied:

\begin{lstlisting}
int c = GetCharPressed();
while (c > 0)
{
    if ((c >= 32) && (c <= 125) && (length < MAX)) name[length++] = (char)c;
    c = GetCharPressed();   // the queue may hold several this frame
}
if (IsKeyPressed(KEY_BACKSPACE) && length > 0) name[--length] = '\0';
\end{lstlisting}

It is a queue, so you loop until it is empty. Handle backspace with
\lstinline|IsKeyPressed| because it is an action, not a character.

\tryit{Lesson 3 shows the pressed-versus-down counters side by side. Hold SPACE
for a second and compare the two numbers.}{./course2d/build.sh 3}
